\chapter{Implementation and Results}

[Must be present in Final Report. Incomplete version might be included in FYDP-1 Report, however it is optional.]

Every chapter should start with 1-2 sentences on the outline of the chapter.

\section{Environment Setup}
\section{Testing and Evaluation}
The Testing and Evaluation section explains how a system is tested and how well it performs, using test cases, metrics, and experimental results to validate its effectiveness. In a thesis (research) work, this typically involves dataset-based evaluation and performance metrics—for example, an AI disease detection model may be tested on unseen data and achieve 91\% accuracy, outperforming an existing model with 85\%, with analysis explaining the improvement. In a project (software/system) work, it focuses more on functional and performance testing—for example, a web-based student management system may be tested with multiple users, showing correct outputs for all test cases and maintaining an average response time of 1.2 seconds under load. In both cases, the section provides measurable evidence, comparison, and analysis to prove that the system meets its objectives


\section{Results and Discussion}
The Results and Discussion section presents the key findings of the work and explains their meaning in relation to the project objectives. It includes a clear presentation of results (tables, graphs, outputs), followed by interpretation—why the results are good or not, what patterns are observed, and how they compare with existing work. In a thesis (research) work, this often involves analyzing model performance and scientific outcomes—for example, an AI model achieving 92\% accuracy may be discussed in terms of improved feature extraction, comparison with earlier studies (e.g., 85\%), and reasons for any limitations such as overfitting or dataset bias. In a project (software/system) work, it focuses on system behavior and usability—for example, a web application showing stable performance with 1.2-second response time under 100 users may be discussed in terms of efficiency, scalability, and user experience, along with any constraints like performance degradation at higher loads. Overall, this section not only reports results but critically interprets them, highlights strengths and weaknesses, and connects outcomes to the original objectives
\section{Summary}